\section{Related Work}

Prior work on rule conflict detection in smart home mainly falls into three lines: policy/property checking, physical-channel-aware analysis, and interaction-pattern-based hybrid detection. Huang et al.~\cite{huang2023survey} note that the problem definition, detection target, and conflict taxonomy are still not fully unified. We compare representative studies by whether they distinguish rule interaction from candidate rule conflicts (DIC) and whether they support user-specific safety preferences with reasonable user effort.

\textbf{Policy/property checking.} Soteria~\cite{celik2018soteria}, IoTGuard~\cite{celik2019iotguard}, SafeChain~\cite{hsu2019safechain}, and IoTSan~\cite{nguyen2018iotsan} verify predefined safety or security properties over IoT apps or trigger--action programs. These approaches are effective for checking explicit properties, but they rely heavily on expert-crafted policies or templates, which are difficult to generalize across diverse smart home deployments and often require substantial manual effort.

\textbf{Physical-channel-aware analysis.} iRuler~\cite{wang2019charting}, IoTIE~\cite{chen2019multi}, IoTSafe~\cite{ding2021iotsafe}, and IoTCom~\cite{alhanahnah2022iotcom} analyze interactions through physical channels to expose hidden dependencies beyond shared devices. However, their support for zone-dependent propagation is limited, absent, or dependent on deployment-specific input, which may over-approximate cross-zone effects and introduce false positives. IoTSafe further uses dynamic monitoring, but at the cost of substantial deployment and testing effort.

\textbf{Interaction-pattern-based hybrid detection.} Pattern-based systems summarize representative rule interaction patterns and flag automation rules that match them~\cite{trimananda2020understanding,chi2020cross,chi2023detecting}. These methods are useful for discovering suspicious rule interaction, but deciding whether a matched pattern should be treated as a candidate rule conflict still requires deployment-specific context. IoTMediator~\cite{chi2023detecting} combines static pre-screening with runtime dynamic monitoring and supports customized handling, yet its pattern-driven target is not directly aligned with DIC under user-specific safety preferences, which can lead to false positives.

Compared with these approaches, our framework uses the Entity Safety Configuration as a lightweight preference baseline to filter rule interaction into candidate rule conflicts, and then applies runtime assertion verification to confirm whether a candidate conflict edge is activated under the current system state. This design eliminates nine false positives reported by IoTMediator in our testbed while keeping the runtime overhead at the millisecond level.